\documentclass[table]{beamer}
%[]中可以使用draft、handout、screen、transparency、trancompress、compress等参数

%指定beamer的模式与主题
\mode<presentation>
{
  \usetheme{Madrid}
%\usetheme{Boadilla}
%\usecolortheme{default}
%\usecolortheme{orchid}
%\usecolortheme{whale}
%\usefonttheme{professionalfonts}
}

%\usetheme{Madrid}
%这里还可以选择别的主题：Bergen, Boadilla, Madrid, AnnArbor, CambridgeUS, Pittsburgh, Rochester, Warsaw, ...
%有导航栏的Antibes, JuanLesPins, Montpellier, ...
%有内容的Berkeley, PaloAlto, Goettingen, Marburg, Hannover, ...
%有最小导航栏的Berlin, Ilmenau, Dresden, Darmstadt, Frankfurt, Singapore, Szeged, ...
%有章和节表单的Copenhagen, Luebeck, Malmoe, Warsaw, ...

%\usecolortheme{default}
%设置内部颜色主题(这些主题一般改变block里的颜色);这个主题一般选择动物来命名
%这里还可以选择别的颜色主题,如默认的和有特别目的的颜色主题default,structure,sidebartab,全颜色主题albatross,beetle,crane,dove,fly,seagull,wolverine,beaver

%\usecolortheme{orchid}
%设置外部颜色主题(这些主题一般改变title里的颜色);这个主题一般选择植物来命名
%这里还可以选择别的颜色主题,如默认的和有特别目的的颜色主题lily,orchid,rose

%\usecolortheme{whale}
%设置字体主题;这个主题一般选择海洋动物来命名
%这里还可以选择别的颜色主题,如默认的和有特别目的的颜色主题whale,seahorse,dolphin

%\usefonttheme{professionalfonts}
%类似的还可以定义structurebold,structuresmallcapsserif,professionalfonts

% 控制 beamer 的风格,可以根据自己的爱好修改
%\usepackage{beamerthemesplit} %使用 split 风格
%\usepackage{beamerthemeshadow} %使用 shadow 风格
%\usepackage[width=2cm,dark,tab]{beamerthemesidebar}

%插入音标
%\usepackage{tipa}
%\AtBeginDocument{
  %\renewcommand\textipa{\fontencoding{T3}\selectfont}
%}
%\AtBeginDocument{
  %\renewcommand\textipa[2][r]{{\fontfamily{cm#1}\tipaencoding #2}}
%}
%\renewenvironment{IPA}[1][r]
 %{\fontfamily{cm#1}\tipaencoding}
 %{}

% 设定英文字体
%\usepackage{fontspec}
% Fix bugs for fontspec in TeXLive2015
\ifdefined\suppressfontnotfounderror
  \expandafter\let\csname xetex_suppressfontnotfounderror:D\endcsname
    \suppressfontnotfounderror
\else
  \expandafter\let\csname xetex_suppressfontnotfounderror:D\endcsname
    \luatexsuppressfontnotfounderror
\fi
\usepackage[no-math]{fontspec}
\setmainfont{Times New Roman}
\setsansfont{Arial}
\setmonofont{Courier New}

% 设定中文字体
\usepackage[BoldFont,SlantFont,CJKchecksingle,CJKnumber]{xeCJK}
%\setCJKmainfont[BoldFont={Adobe Heiti Std},ItalicFont={Adobe Kaiti Std}]{Adobe Song Std}
\setCJKmainfont[BoldFont={Adobe Heiti Std},ItalicFont={Adobe Kaiti Std}]{WenQuanYi Micro Hei}
\setCJKsansfont{Adobe Heiti Std}
\setCJKmonofont{Adobe Fangsong Std}
\punctstyle{hangmobanjiao}

\defaultfontfeatures{Mapping=tex-text}
\usepackage{xunicode}
\usepackage{xltxtra}

\XeTeXlinebreaklocale "zh"
\XeTeXlinebreakskip = 0pt plus 1pt minus 0.1pt

\usepackage{setspace}
\usepackage{colortbl,xcolor}
\usepackage{hyperref}
%\hypersetup{xetex,bookmarksnumbered=true,bookmarksopen=true,pdfborder=1,breaklinks,colorlinks,linkcolor=blue,filecolor=black,urlcolor=cyan,citecolor=green}
\hypersetup{xetex,bookmarksnumbered=true,bookmarksopen=true,pdfborder=1,breaklinks,colorlinks,linkcolor=cyan,filecolor=black,urlcolor=blue,citecolor=green}

% 插入图片
\usepackage{graphicx}
\graphicspath{{figures/}}
% 图文混排
%\usepackage{picins}
\usepackage{floatflt}

% 可能用到的包
\usepackage{amsmath,amssymb}
%插入多媒体
%\usepackage{media9}
%\usepackage{movie15}
\usepackage{multimedia}
\usepackage{multicol}
\usepackage{multirow}

% 定义一些自选的模板，包括背景、图标、导航条和页脚等，修改要慎重
% 设置背景渐变由10%的红变成10%的结构颜色
%\beamertemplateshadingbackground{red!10}{structure!10}
%\beamertemplatesolidbackgroundcolor{white!90!blue}
% 使所有隐藏的文本完全透明、动态，而且动态的范围很小
\beamertemplatetransparentcovereddynamic
% 使itemize环境中变成小球，这是一种视觉效果
\beamertemplateballitem
% 为所有已编号的部分设置一个章节目录，并且编号显示成小球
\beamertemplatenumberedballsectiontoc
% 将每一页的要素的要素名设成加粗字体
\beamertemplateboldpartpage

% item逐步显示时,使已经出现的item、正在显示的item、将要出现的item呈现不同颜色
\def\hilite<#1>{
 \temporal<#1>{\color{gray}}{\color{blue}}
    {\color{blue!25}}
}

\renewcommand{\today}{\number\year 年 \number\month 月 \number\day 日}

%五角星
\usepackage{MnSymbol}

%去除图表标题中的figure等
\usepackage{caption}
\captionsetup{labelformat=empty,labelsep=none}

\usepackage{tabu}
\usepackage{multirow}
%表格自动换行
\usepackage{tabularx} 

% 千分号
%\usepackage{textcomp}

%罗马数字
\makeatletter
\newcommand{\rmnum}[1]{\romannumeral #1}
\newcommand{\Rmnum}[1]{\expandafter\@slowromancap\romannumeral #1@}
\makeatother

%分栏
\usepackage{multicol}

%\usepackage{enumitem}
%\usepackage{enumerate}

%键盘
\usepackage{keystroke}

%心形
%\usepackage{fdsymbol}

%插入源代码
\usepackage{listings}
\lstset{
  language=perl,                  % 程序语言名称：TeX, Perl, R, sh, bash, Awk
  basicstyle=\normalsize\tt,      %\tt指monospace字体族，程序源代码使用此族字体表示更加美观
  numbers=left,                   % 行号位置(左侧)
  numberstyle=\small,             % 行号字体的字号
  stepnumber=1,                   % 行号的显示步长
  numbersep=5pt,                  % 行号与代码间距
  backgroundcolor=\color{white},  % 背景色;需要 \usepackage{color}
  showspaces=false,               % 不显示空格
  showstringspaces=false,         % 不显示代码字符串中的空格标记
  showtabs=false,                 % 不显示 TAB
  tabsize=4, 
  frame=shadowbox,                % 把代码用带有阴影的框圈起来
  captionpos=b,                   % 标题位置
  breaklines=true,                % 对过长的代码自动断行
  breakatwhitespace=false,        % 断行只在空格处
  extendedchars=false,            % 解决代码跨页时，章节标题，页眉等汉字不显示的问题
  %escapeinside={\%*}{*},         % 跳脱字符，添加注释，暂时离开 listings 
  %escapeinside=``,
  commentstyle=\color{red!50!green!50!blue!50}\tt,  %浅灰色的注释
  rulesepcolor=\color{red!20!green!20!blue!20},     %代码块边框为淡青色
  keywordstyle=\color{blue!70}\bfseries\tt,         %代码关键字的颜色为蓝色，粗体
  identifierstyle=\tt,
  stringstyle=\tt,                % 代码字符串的特殊格式
  keepspaces=true,
  breakindent=1em,
  %breakindent=22pt,
  %breakindent=4em,
  breakautoindent=true,
  flexiblecolumns=true,
  aboveskip=1em,                  %代码块边框
  xleftmargin=2em,
  xrightmargin=2em
}

%\setbeamercolor{alerted text}{fg=magenta}
\setbeamercolor{bgcolor}{fg=yellow,bg=cyan}
%\setbeamercolor{itemize/enumerate body}{fg=green}

\begin{document}

%\includeonlyframes{current}

\logo{\includegraphics[height=0.08\textwidth]{qr.png}}

% 在每个Section前都会加入的Frame
\AtBeginSection[]
{
  \begin{frame}<beamer>
    %\frametitle{Outline}
    \frametitle{教学提纲}
    \setcounter{tocdepth}{3}
    \begin{multicols}{2}
      \tableofcontents[currentsection,currentsubsection]
      %\tableofcontents[currentsection]
    \end{multicols}
  \end{frame}
}
% 在每个Subsection前都会加入的Frame
\AtBeginSubsection[]
{
  \begin{frame}<beamer>
%%\begin{frame}<handout:0>
%% handout:0 表示只在手稿中出现
    \frametitle{教学提纲}
    \setcounter{tocdepth}{3}
    \begin{multicols}{2}
    \tableofcontents[currentsection,currentsubsection]
    \end{multicols}
%% 显示在目录中加亮的当前章节
  \end{frame}
}

% 为当前幻灯片设置背景
%{
%\usebackgroundtemplate{
%\vbox to \paperheight{\vfil\hbox to
%\paperwidth{\hfil\includegraphics[width=2in]{tijmu_charcoal.png}\hfil}\vfil}
%}
\begin{frame}[plain]
  \begin{center}
    {\Huge 疯狂的实验\\}
    \vspace{1cm}
    {\LARGE 天津医科大学\\}
    %\vspace{0.2cm}
    {\LARGE 生物医学工程与技术学院\\}
    \vspace{1cm}
    {\large 2018-2019学年下学期（春）\\ 公共选修课}
  \end{center}
\end{frame}
%}



%\includeonlyframes{current}

\title[生物学]{第四章\quad 疯狂的实验——生物学}
\author[Yixf]{伊现富（Yi Xianfu）}
\institute[TIJMU]{天津医科大学（TIJMU）\\ 生物医学工程与技术学院}
\date{2017年3月}

\begin{frame}
  \titlepage
\end{frame}

\begin{frame}[plain,label=current]
  \frametitle{教学提纲}
  \setcounter{tocdepth}{3}
  \begin{multicols}{2}
    \tableofcontents
  \end{multicols}
\end{frame}


% 
\section{}

% 
\section{}

% 
\section{}

% 
\section{}

% 
\section{}


\section{史海撷华}
% 1-6,1620
\begin{frame}
  \frametitle{史海撷华 | 光合作用}
  \begin{block}{种柳行动}
    \begin{itemize}
      \item 范\textbullet 黑尔蒙特（最后一个炼金术士，第一个化学家），首位用泥土、树木和称实际操作实验的人
      \item 200磅在炉中烘干的泥，5磅重的柳树幼枝，定期浇水
      \item 5年后拔出柳树，对土和柳树分别称重
    \end{itemize}
  \end{block}
  \vspace{-0.75em}
  \pause
  \begin{block}{结果与结论}
    \begin{itemize}
      \item 泥土减重2盎司，树木增重164磅零3盎司
      \item 164磅的木质、树皮以及根系都只来源于水（当时唯一合理的结论）
    \end{itemize}
  \end{block}
  \vspace{-0.75em}
  \pause
  \begin{block}{启示}
    \begin{itemize}
      \item 为实验铺平了道路，使其从此成为获取认识的手段
      \item 他的想法启发了很多的学者，开展罐中植物的研究
      \item 为“光合作用”这一神秘过程的探究开了先河
      \item 学生借此测试洞察力，练习严谨的实验设计
    \end{itemize}
  \end{block}
\end{frame}

% 1-8,1729
\begin{frame}
  \frametitle{史海撷华 | 生物钟 | 植物}
  \begin{block}{现象}
    \begin{itemize}
      \item 含羞草在夜间合拢叶片，白天打开。
      \item 如果把含羞草置于一个它无法知晓昼夜的环境中，情况会怎样呢？
    \end{itemize}
  \end{block}
  \pause
  \begin{block}{实验}
    \begin{itemize}
      \item 1729年，让\textbullet 雅克\textbullet 徳奥图斯\textbullet 德迈朗（天文学家）
      \item 把一株含羞草放到漆黑的柜中
    \end{itemize}
  \end{block}
  \pause
  \begin{block}{结果}
    \begin{itemize}
      \item 叶片在没有太阳光的情况下，还可以有规律地开合。
    \end{itemize}
  \end{block}
  \pause
  \begin{block}{启示}
    \begin{itemize}
      \item “时间生物学”（研究生物体内部的生物钟）的创立者
    \end{itemize}
  \end{block}
\end{frame}

% 1-109,1938
\begin{frame}
  \frametitle{史海撷华 | 生物钟 | 人}
  \begin{block}{疑问}
    人的睡眠规律究竟只是习惯，抑或是人体内存在着生物钟？
  \end{block}
  \pause
  \begin{block}{实验}
    \begin{itemize}
      \item 实验一：把生物钟从每天24小时调整为48小时、12小时——实验无果而终
      \item 实验二：把生物钟从每天24小时调整为21小时、28小时，测量体温——结果模棱两可
      \item 实验三：宽20米、高8米的猛犸洞窟（漆黑、安静、恒温）——结果显示出两面性
    \end{itemize}
  \end{block}
  \pause
  \begin{block}{启示}
    随后的实验证实，人体内确实存在着生物钟。它的运转大致跟一天24小时相吻合，并且每天都会根据实际时长进行自动调整。
  \end{block}
\end{frame}

% 1-12,1774
\begin{frame}
  \frametitle{史海撷华 | 温度耐受}
  \begin{block}{疑问}
    人体能够承受什么样的温度？
  \end{block}
  \pause
  \begin{block}{实验}
    \begin{itemize}
      \item 建造桑拿室，45-100-127
      \item 穿衣服发汗，赤裸着手持一只平底锅，上面放着一块牛排
    \end{itemize}
  \end{block}
  \pause
  \begin{block}{结果及结论}
    \begin{itemize}
      \item 24页的《皇家协会学报》
      \item 有一个“与生命体直接相关的自然的系统”消除热量？（降温——汗液等的蒸发加之以血液流动！）
    \end{itemize}
  \end{block}
\end{frame}

% 1-16,1802
\begin{frame}
  \frametitle{史海撷华 | 蛙腿实验}
  \begin{block}{现象}
    如果用两种不同的金属触碰蛙腿，连成通路，它们的肌肉会抽搐。
  \end{block}
  \pause
  \begin{block}{理论}
    \begin{itemize}
      \item 加尔瓦尼：“动物电流”与生命力有关，效果与电流通过无生命的物质是不同的。
      \item 伏特：世界上只有一种电，如论市雷雨天的闪电还是抽搐的蛙腿，原理都与这种电有关。
    \end{itemize}
  \end{block}
  \pause
  \begin{block}{启示}
    \begin{itemize}
      \item 《弗兰肯斯坦》（第一部科幻小说）——玛丽\textbullet 雪莱
      \item 眨眼睛的尸体——对于绞死者进行头颅实验
    \end{itemize}
  \end{block}
\end{frame}

% 1-20,1825
\begin{frame}
  \frametitle{史海撷华 | 胃上有洞的人}
  \begin{block}{实验材料}
       1882年，威廉\textbullet 博蒙特；亚力克西斯\textbullet 圣马丁（受伤的士兵）
  \end{block}
  \vspace{-0.5em}
  \pause
  \begin{block}{疑问}
消化仅仅是一个纯化学的过程，还是同时需要人体提供某种未知的生命力量促使其完成？消化和腐烂的区别是否在于前者拥有人体内的未知生命力量，而后者没有？
  \end{block}
  \vspace{-0.5em}
  \pause
  \begin{block}{实验}
      用丝线系好的食物，塞入胃中——拉出，观察消化情况；插入软管、导出胃液，把一把牛肉粒置于其中；……
  \end{block}
  \vspace{-0.5em}
  \pause
  \begin{block}{结果}
    \begin{itemize}
      \item 器皿中胃液的化学能量足以完成消化（无需未知的生命力量）。
      \item 推翻了胃液仅仅是流于胃中储存起来的唾液的推测。
      \item 医学伦理
    \end{itemize}
  \end{block}
\end{frame}

% 
\begin{frame}
  \frametitle{史海撷华 | }
\end{frame}

% 
\begin{frame}
  \frametitle{史海撷华 | }
\end{frame}

% 
\begin{frame}
  \frametitle{史海撷华 | }
\end{frame}

% 
\begin{frame}
  \frametitle{史海撷华 | }
\end{frame}

% 
\begin{frame}
  \frametitle{史海撷华 | 其他}
  \begin{block}{真正的疯狂}
    \begin{itemize}
      \item 蚯蚓没有听觉——达尔文为蚯蚓演奏巴松管、笛子和钢琴
      \item 
      \item 
      \item 
      \item 
      \item 
      \item 
      \item 
      \item 
    \end{itemize}
  \end{block}
\end{frame}



\section*{Acknowledgements}
\begin{frame}
  \frametitle{Powered by}
  \begin{center}
    \includegraphics[width=9cm]{power.png}
  \end{center}
\end{frame}

\end{document}

